\documentclass{beamer}

\usepackage[utf8]{inputenc}
\usepackage[T1]{fontenc}
\usepackage[spanish]{babel}
\usepackage{amsmath, amssymb}
\usepackage{graphicx}
\usepackage{booktabs}
\usepackage{tikz}
\usepackage{pgfplots}
\pgfplotsset{compat=1.18}

\usetheme{Madrid}

\title{Economía Poblacional}
\subtitle{México y América Latina: envejecimiento, economía y política pública}
\author{Héctor J. Villarreal}
\date{23 de marzo de 2026}

\begin{document}

%------------------------------------------------
\begin{frame}
\titlepage
\end{frame}

%------------------------------------------------
\begin{frame}{Objetivo de esta segunda presentación}
\begin{itemize}
    \item Aplicar al caso de América Latina y México la discusión del curso sobre envejecimiento.
    \item Sintetizar los principales hallazgos recientes de CEPAL.
    \item Discutir los efectos económicos del envejecimiento sobre:
    \begin{itemize}
        \item crecimiento,
        \item mercado laboral,
        \item consumo,
        \item finanzas públicas,
        \item y estructura sectorial.
    \end{itemize}
    \item Extraer implicaciones para México y para una agenda de política pública.
\end{itemize}
\end{frame}

%------------------------------------------------
\begin{frame}{Mensaje central}
\begin{itemize}
    \item América Latina envejece con rapidez.
    \item El proceso ocurre con menor ingreso per cápita, mayor desigualdad y Estados de bienestar menos consolidados que en Europa.
    \item El envejecimiento es simultáneamente:
    \begin{itemize}
        \item un logro social,
        \item un desafío macrofiscal,
        \item y una oportunidad de transformación productiva.
    \end{itemize}
\end{itemize}
\end{frame}

%------------------------------------------------
\begin{frame}{Un punto clave para el curso}
\begin{itemize}
    \item No estamos frente a un cambio marginal.
    \item Estamos frente a una transformación estructural de la economía:
    \begin{itemize}
        \item cambia la oferta de trabajo,
        \item cambia la composición del consumo,
        \item cambia la dirección de las transferencias intergeneracionales,
        \item y cambia la presión sobre el Estado.
    \end{itemize}
    \item Esto es particularmente importante en América Latina.
\end{itemize}
\end{frame}

%------------------------------------------------
\begin{frame}{Magnitud del envejecimiento en América Latina y el Caribe}
\begin{itemize}
    \item En 2024, la población de 65 años y más se estima en alrededor de 65 millones.
    \item Eso equivale a 9.9\% de la población total.
    \item Para 2050, este grupo alcanzaría 138 millones.
    \item Su peso relativo subiría a 18.9\% del total regional.
\end{itemize}

\vspace{0.3cm}
\textbf{Lectura:} la proporción prácticamente se duplica en poco más de un cuarto de siglo.
\end{frame}

%------------------------------------------------
\begin{frame}{Gráfica 1: el cambio etario de la región}
\centering
\begin{tikzpicture}
\begin{axis}[
    width=10.8cm,
    height=5.9cm,
    xlabel={Año},
    ylabel={Porcentaje de la población},
    xmin=2000, xmax=2100,
    ymin=0, ymax=80,
    legend style={at={(0.98,0.98)},anchor=north east,font=\scriptsize},
    xtick={2000,2025,2050,2075,2100},
    ytick={0,20,40,60,80}
]
\addplot[thick] coordinates {(2000,32) (2025,25) (2050,18) (2075,15) (2100,14)};
\addlegendentry{0-14}
\addplot[thick,dashed] coordinates {(2000,63) (2025,65) (2050,61) (2075,55) (2100,49)};
\addlegendentry{15-64}
\addplot[thick,dotted] coordinates {(2000,5) (2025,10) (2050,19) (2075,28) (2100,32)};
\addlegendentry{65+}
\end{axis}
\end{tikzpicture}

\vspace{0.2cm}
\small
Perfil estilizado a partir del diagnóstico de CEPAL: cae la población joven, primero domina la edad laboral, y luego aumenta con fuerza la población mayor.
\end{frame}

%------------------------------------------------
\begin{frame}{Más rápido de lo que esperábamos}
\begin{itemize}
    \item En 1950, la región tenía una proporción de población de 65 años y más de 3.2\%.
    \item Hacia 2100, CEPAL reporta una proyección de 31.7\%.
    \item La región duplicará su actual proporción de población de 65 años y más en apenas 28 años.
    \item Llegaría a cerca de 19.7\% en 2052.
\end{itemize}

\vspace{0.3cm}
\textbf{Conclusión:} el envejecimiento avanza más rápido que la capacidad de ajuste institucional.
\end{frame}

%------------------------------------------------
\begin{frame}{Envejecimiento del envejecimiento}
\begin{itemize}
    \item No sólo aumenta la población mayor.
    \item Crece con especial fuerza la población de 80 años y más.
    \item Según CEPAL:
    \begin{itemize}
        \item pasa de 12 millones en 2024,
        \item a 37 millones en 2050.
    \end{itemize}
    \item Su peso relativo subiría de 1.9\% a 5.1\%.
\end{itemize}

\vspace{0.3cm}
Esto importa porque en ese grupo aumenta con más fuerza la probabilidad de dependencia funcional y de demanda de cuidados de largo plazo.
\end{frame}

%------------------------------------------------
\begin{frame}{Cuatro rasgos que vuelven más difícil el ajuste en LATAM}
\begin{enumerate}
    \item Menor nivel de desarrollo económico.
    \item Ritmo de envejecimiento más acelerado.
    \item Estados de bienestar menos consolidados.
    \item Desigualdad estructural más alta.
\end{enumerate}

\vspace{0.3cm}
Por eso América Latina no enfrenta simplemente el mismo problema que Europa, pero con menos años de retraso: enfrenta una variante más dura.
\end{frame}

%------------------------------------------------
\begin{frame}{Dependencia demográfica: intuición básica}
\[
\text{Relación de dependencia}=
\frac{\text{Población menor de 15}+\text{Población mayor de 65}}{\text{Población de 15 a 64}}
\]

\begin{itemize}
    \item Primero baja por la reducción de la dependencia infantil.
    \item Luego vuelve a subir por el crecimiento de la dependencia de personas mayores.
\end{itemize}

\vspace{0.3cm}
El cambio en esta relación resume una parte importante del trasfondo económico del envejecimiento.
\end{frame}

%------------------------------------------------
\begin{frame}{Gráfica 2: fin del bono demográfico}
\centering
\begin{tikzpicture}
\begin{axis}[
    width=10.8cm,
    height=5.9cm,
    xlabel={Año},
    ylabel={Dependientes por 100 personas de 15-64},
    xmin=1950, xmax=2050,
    ymin=0, ymax=100,
    legend style={at={(0.98,0.98)},anchor=north east,font=\scriptsize},
    xtick={1950,1975,2000,2025,2050},
    ytick={0,20,40,60,80,100}
]
\addplot[thick] coordinates {(1950,80) (1970,73) (1990,62) (2010,53) (2025,55) (2050,68)};
\addlegendentry{Dependencia total}
\addplot[thick,dashed] coordinates {(1950,72) (1970,63) (1990,50) (2010,38) (2025,31) (2050,26)};
\addlegendentry{Dependencia infantil}
\addplot[thick,dotted] coordinates {(1950,8) (1970,10) (1990,12) (2010,15) (2025,24) (2050,42)};
\addlegendentry{Dependencia de mayores}
\end{axis}
\end{tikzpicture}

\vspace{0.2cm}
\small
Idea central del gráfico de CEPAL: la baja de la dependencia infantil ya no compensa el aumento de la dependencia de personas mayores.
\end{frame}

%------------------------------------------------
\begin{frame}{El bono demográfico está por terminar}
\begin{itemize}
    \item CEPAL proyecta que el bono demográfico regional concluye, en promedio, en 2028.
    \item Existen grandes diferencias entre países, pero el diagnóstico regional es claro:
    \begin{itemize}
        \item la ventana se está cerrando,
        \item y en varios países ya se cerró.
    \end{itemize}
\end{itemize}

\vspace{0.3cm}
El problema no es sólo demográfico: es que la región no aprovechó el bono tan bien como Asia oriental.
\end{frame}

%------------------------------------------------
\begin{frame}{Bono desaprovechado}
\begin{itemize}
    \item CEPAL subraya que América Latina no transformó su bono demográfico en ganancias económicas comparables a Asia oriental.
    \item El dividendo no era automático.
    \item Requería:
    \begin{itemize}
        \item creación de empleo,
        \item inversión en capacidades humanas,
        \item mayor participación de las mujeres,
        \item y un entorno favorable a la productividad.
    \end{itemize}
\end{itemize}
\end{frame}

%------------------------------------------------
\begin{frame}{Mercado laboral: más personas mayores trabajando}
\begin{itemize}
    \item En 2024, aproximadamente un cuarto de las personas de 65 años y más participa en el mercado laboral en América Latina.
    \item La participación es mayor entre hombres.
    \item Esto refleja dos fuerzas:
    \begin{itemize}
        \item deseo de seguir activos,
        \item necesidad de complementar ingresos.
    \end{itemize}
\end{itemize}

\vspace{0.3cm}
En una región con pensiones insuficientes, trabajo y vejez se vuelven temas inseparables.
\end{frame}

%------------------------------------------------
\begin{frame}{Pensiones insuficientes y participación laboral}
\begin{itemize}
    \item CEPAL encuentra una correlación positiva entre:
    \begin{itemize}
        \item insuficiencia de pensiones,
        \item y participación laboral de personas mayores.
    \end{itemize}
    \item Para muchas personas mayores, trabajar no es sólo opción: es necesidad.
\end{itemize}

\vspace{0.3cm}
Esto es crucial para México, donde la informalidad complica la construcción de trayectorias contributivas suficientes.
\end{frame}

%------------------------------------------------
\begin{frame}{Otra palanca clave: la participación de las mujeres}
\begin{itemize}
    \item CEPAL insiste en que una parte del ajuste puede venir de una mayor participación laboral femenina.
    \item Esto tiene doble relevancia:
    \begin{itemize}
        \item compensa parcialmente la desaceleración de la fuerza de trabajo,
        \item y mejora el ingreso agregado de los hogares.
    \end{itemize}
    \item Pero exige políticas de cuidado, inclusión laboral y reducción de desigualdades.
\end{itemize}
\end{frame}

%------------------------------------------------
\begin{frame}{NTA: el lenguaje adecuado para pensar el problema}
\[
C(a)-Y_L(a)=\tau(a)+\big[Y_A(a)-S(a)\big]
\]

\begin{itemize}
    \item $C(a)$: consumo a la edad $a$.
    \item $Y_L(a)$: ingreso laboral.
    \item $\tau(a)$: transferencias netas.
    \item $Y_A(a)-S(a)$: reasignación basada en activos.
\end{itemize}

\vspace{0.3cm}
Las Cuentas Nacionales de Transferencias permiten ver cómo se financia el ciclo de vida.
\end{frame}

%------------------------------------------------
\begin{frame}{Tres etapas del ciclo de vida económico}
\begin{itemize}
    \item Primera etapa: infancia y juventud
    \begin{itemize}
        \item consumo mayor al ingreso laboral.
    \end{itemize}
    \item Segunda etapa: edades laborales
    \begin{itemize}
        \item ingreso laboral mayor al consumo.
    \end{itemize}
    \item Tercera etapa: vejez
    \begin{itemize}
        \item vuelve a aparecer un déficit.
    \end{itemize}
\end{itemize}

\vspace{0.3cm}
El problema económico es cómo se financian esas etapas deficitarias.
\end{frame}

%------------------------------------------------
\begin{frame}{Gráfica 3: perfil estilizado de ingreso y consumo}
\centering
\begin{tikzpicture}
\begin{axis}[
    width=10.8cm,
    height=5.9cm,
    xlabel={Edad},
    ylabel={Índice},
    xmin=0, xmax=90,
    ymin=0, ymax=1.2,
    legend style={at={(0.98,0.98)},anchor=north east,font=\scriptsize},
    xtick={0,20,40,60,80},
    ytick={0,0.2,0.4,0.6,0.8,1.0,1.2}
]
\addplot[thick] coordinates {(0,0.38) (10,0.55) (20,0.72) (30,0.90) (40,0.84) (50,0.86) (60,0.88) (70,0.86) (80,0.82) (90,0.90)};
\addlegendentry{Consumo}
\addplot[thick,dashed] coordinates {(0,0.00) (10,0.00) (20,0.15) (25,0.60) (30,0.90) (35,1.05) (45,1.08) (55,1.02) (60,0.88) (65,0.55) (70,0.30) (80,0.08) (90,0.00)};
\addlegendentry{Ingreso laboral}
\end{axis}
\end{tikzpicture}

\vspace{0.2cm}
\small
El perfil sintetiza el patrón mostrado por CEPAL para seis países: dos etapas deficitarias y una superavitaria.
\end{frame}

%------------------------------------------------
\begin{frame}{La gran transformación que viene}
\begin{itemize}
    \item Hoy, en la región, el déficit agregado de la infancia todavía domina.
    \item Pero CEPAL proyecta que, con el tiempo, el déficit de la vejez ganará peso.
    \item Las necesidades de financiamiento se desplazarán progresivamente hacia el final del ciclo de vida.
\end{itemize}

\vspace{0.3cm}
Este punto es central para pensar pensiones, salud y cuidados.
\end{frame}

%------------------------------------------------
\begin{frame}{Segundo bono demográfico}
\begin{itemize}
    \item CEPAL recupera la idea del segundo bono demográfico.
    \item La lógica es:
    \begin{itemize}
        \item poblaciones más longevas acumulan más activos para financiar una vejez más larga,
        \item si esos activos se invierten productivamente, puede aumentar la productividad.
    \end{itemize}
    \item No es automático.
\end{itemize}

\vspace{0.3cm}
Requiere ahorro, instituciones y mercados de capital más profundos.
\end{frame}

%------------------------------------------------
\begin{frame}{Pero América Latina enfrenta límites}
\begin{itemize}
    \item Menor capacidad de ahorro.
    \item Mercados financieros menos desarrollados.
    \item Cobertura previsional incompleta.
    \item Alta informalidad.
\end{itemize}

\vspace{0.3cm}
Por eso el segundo bono demográfico es una posibilidad, no una garantía.
\end{frame}

%------------------------------------------------
\begin{frame}{PIB per cápita: una descomposición útil}
\[
\frac{Y}{P}=\frac{Y}{L}\times\frac{L}{PET}\times\frac{PET}{P}
\]

\begin{itemize}
    \item $\frac{Y}{L}$: productividad laboral.
    \item $\frac{L}{PET}$: tasa de empleo.
    \item $\frac{PET}{P}$: proporción de población en edad de trabajar.
\end{itemize}

\vspace{0.3cm}
El envejecimiento afecta directamente los dos últimos términos, y por eso presiona el PIB per cápita.
\end{frame}

%------------------------------------------------
\begin{frame}{Lo más importante: productividad}
\begin{itemize}
    \item CEPAL es clara:
    \begin{itemize}
        \item el impacto demográfico negativo no es un destino inexorable,
        \item puede mitigarse con productividad.
    \end{itemize}
    \item También ayudan:
    \begin{itemize}
        \item mayor participación laboral femenina,
        \item mejor salud,
        \item vidas laborales más largas,
        \item e innovación.
    \end{itemize}
\end{itemize}
\end{frame}

%------------------------------------------------
\begin{frame}{La heterogeneidad regional importa}
\begin{itemize}
    \item CEPAL muestra que el impacto del envejecimiento sobre el PIB per cápita entre 2025 y 2050 no será idéntico en todos los países.
    \item Algunos ya enfrentarán presión negativa clara.
    \item Otros todavía conservan algo de impulso demográfico.
\end{itemize}

\vspace{0.3cm}
Eso obliga a pensar una agenda regional, pero también agendas nacionales diferenciadas.
\end{frame}

%------------------------------------------------
\begin{frame}{Pobreza y vulnerabilidad en la vejez}
\begin{itemize}
    \item En la región, la tasa de pobreza de 65 años y más es menor que la de niños y jóvenes.
    \item CEPAL reporta cerca de 14.7\% para personas de 65 años y más.
    \item Entre menores de 15 años, la tasa ronda 41.4\%.
\end{itemize}

\vspace{0.3cm}
Pero esto no significa ausencia de vulnerabilidad: significa que la vejez depende críticamente de redes de apoyo, patrimonio y protección social.
\end{frame}

%------------------------------------------------
\begin{frame}{Desigualdad y vejez}
\begin{itemize}
    \item La vulnerabilidad no está distribuida homogéneamente.
    \item CEPAL subraya mayor exposición para:
    \begin{itemize}
        \item mujeres,
        \item pueblos indígenas,
        \item afrodescendientes,
        \item migrantes,
        \item personas con discapacidad.
    \end{itemize}
    \item En otras palabras, la vejez hereda desigualdades acumuladas durante la vida.
\end{itemize}
\end{frame}

%------------------------------------------------
\begin{frame}{Patrones de consumo: cambia la composición, no sólo el nivel}
\begin{itemize}
    \item Los hogares con personas mayores muestran patrones de gasto distintos.
    \item Mayor peso relativo de:
    \begin{itemize}
        \item alojamiento,
        \item salud,
        \item cuidado,
        \item servicios personales.
    \end{itemize}
    \item Menor peso relativo de educación.
\end{itemize}

\vspace{0.3cm}
El envejecimiento recompone la demanda agregada.
\end{frame}

%------------------------------------------------
\begin{frame}{Economía plateada: una oportunidad real, pero no automática}
\begin{itemize}
    \item CEPAL discute tres marcos:
    \begin{itemize}
        \item economía plateada,
        \item economía de la longevidad,
        \item economía generacional.
    \end{itemize}
    \item El punto común:
    \begin{itemize}
        \item una sociedad más longeva genera nuevas demandas,
        \item y esas demandas pueden impulsar sectores productivos.
    \end{itemize}
\end{itemize}
\end{frame}

%------------------------------------------------
\begin{frame}{Sectores con mayor potencial en una región que envejece}
\begin{itemize}
    \item Salud.
    \item Cuidados y cuidados de largo plazo.
    \item Industria farmacéutica y biotecnología.
    \item Sector financiero y asegurador.
    \item Tecnología y servicios digitales adaptados.
    \item Turismo senior.
    \item Vivienda y urbanismo adaptado.
\end{itemize}

\vspace{0.3cm}
La clave es que el envejecimiento no sólo presiona gasto: también redirige inversión y empleo.
\end{frame}

%------------------------------------------------
\begin{frame}{¿Qué significa esto para México?}
\begin{itemize}
    \item México comparte los grandes rasgos del diagnóstico regional:
    \begin{itemize}
        \item envejecimiento acelerado,
        \item informalidad alta,
        \item pensiones desiguales,
        \item presión futura en salud y cuidados.
    \end{itemize}
    \item Esto vuelve especialmente importante:
    \begin{itemize}
        \item la productividad,
        \item la formalización,
        \item y la arquitectura fiscal.
    \end{itemize}
\end{itemize}
\end{frame}

%------------------------------------------------
\begin{frame}{Lectura DFD del problema}
\begin{itemize}
    \item El envejecimiento no es sólo un tema de bienestar.
    \item Es también un problema de hoja de balance intertemporal:
    \begin{itemize}
        \item más pasivos implícitos en pensiones,
        \item más gasto en salud,
        \item más necesidad de cuidados,
        \item menor impulso demográfico al crecimiento.
    \end{itemize}
    \item Por eso conecta directamente con deuda y sostenibilidad.
\end{itemize}
\end{frame}

%------------------------------------------------
\begin{frame}{Cierre final}
\begin{itemize}
    \item América Latina envejece rápido.
    \item México no está fuera de esa historia.
    \item El envejecimiento trae:
    \begin{itemize}
        \item desafíos macroeconómicos,
        \item presiones fiscales,
        \item y oportunidades productivas.
    \end{itemize}
    \item La diferencia entre crisis y adaptación dependerá de la calidad de la política pública.
\end{itemize}

\vspace{0.5cm}
\centering
\Large
\textbf{La región ya no puede pensar crecimiento sin pensar demografía.}
\end{frame}

\end{document}